% Estrutura sugerida para a introdução
%%%%%%%%%%%%%%%%%%%%%%%%%%%%%%%%%%%%%%%
% (1) Contexto geral do problema
%     Apresente o tema principal da tese e o contexto matemático;

% (2) Motivação
%     Explique por que o problema é relevante na literatura.
%     Cite trabalhos clássicos e resultados conhecidos;
%
% (3) Problema central da tese
%     Descreva claramente qual é a questão investigada ou qual lacuna
%     na literatura o trabalho pretende abordar;
%
% (4) Principais resultados
%     Enuncie de forma informal os teoremas principais obtidos;
%
% (5) Ideia das demonstrações
%     Dê uma breve visão das técnicas utilizadas;
%
% (6) Organização da tese
%     Descreva brevemente o conteúdo de cada capítulo.
%%%%%%%%%%%%%%%%%%%%%%%%%%%%%%%%%%%%%%%%%%%%%%%%%%%%%%

\chapter*{Introdução}
\addcontentsline{toc}{chapter}{Introdução}

% Comece aqui!